\documentclass[10pt,a4paper]{amsart}
\usepackage[margin=19mm]{geometry}
\usepackage{amsmath,amssymb,mathtools}
\usepackage[scr=rsfs,cal=euler]{mathalfa}
\usepackage{xfrac}
\usepackage[unicode,hidelinks,bookmarks=false]{hyperref}
\newcommand{\ie}{i.e.\ }
\newcommand{\cf}{cf.\ }
\newcommand{\resp}{respectively\ }
\newcommand{\Subsection}{\subsection}
\providecommand{\nobreakdash}{\nobreak\hbox{-}}
\setlength{\emergencystretch}{2em}
\multlinegap0pt
\usepackage{algorithm}\usepackage{algpseudocode}
\usepackage[shortlabels]{enumitem}
\usepackage{tikz}
\newtheorem{definition}{Definition}
\newtheorem{lemma}{Lemma}
\title{Temporal Role Colouring}
\author{Jessica Enright, Kitty Meeks, Puck Rombach, Ella Yates}
\date{Source: arXiv:2607.29272v1 (CC BY 4.0)}
\begin{document}
\maketitle

\noindent\emph{Source and licence.} Temporal Role Colouring. Version arXiv:2607.29272v1 (CC BY 4.0), distributed by arXiv under the Creative Commons Attribution 4.0 International licence, http://creativecommons.org/licenses/by/4.0/, as recorded in the arXiv licence metadata for this exact version and shown on the abstract page https://arxiv.org/abs/2607.29272v1. Full licence notice: \texttt{licenses/arxiv-2607.29272-CC-BY-4.0.txt}.\par\smallskip

\noindent\textit{Editorial scope.} Counted unit: the article's lemma state (source0.tex lines 513--515). The article's complete original proof of it is delivered with it (source0.tex lines 516--537, one \texttt{proof} environment of 22 lines). No mathematics is written, repaired or re-derived: the statement and the proof are the article's own lines, in the article's own notation, and no retained line is substituted or re-flowed.  Retained uncounted as the source-local context the proof needs: lines 486--506.  Nothing was omitted from any retained span, so the package carries no omission record.  No retained line contains a citation, so the wrapper carries no bibliography and no bibliography entry.  No retained line contains a figure environment, an image-inclusion command or drawing code, so no image is delivered and the package declares no asset dependency.  Editorial formatting only, in addition: an AMS \texttt{amsart} preamble that restates the article's own declarations and macros used by the retained lines, namely the environments \texttt{definition}, \texttt{lemma} on separate counters, together with the standard packages those lines need.  The retained environments print in this document as Definition 1, Lemma 1 in order of appearance, whereas the article prints the same items with the numbering of its own full document.  The complete native source of the article as distributed in the arXiv e-print for this exact version is bundled unchanged as \texttt{sources/206469/source0.tex}.
\begin{definition}[(k,X,f)-Component-Exchangeable Temporally Uniform Problem] \label{def:kxf problem}
We say that a decision problem $P$ with input $x = (\mathcal{G},\beta)$ such that $\mathcal{G}$ has lifetime $\Lambda$ is \emph{$(k,X,f)$
component-exchangeable temporally uniform} if and only if there exist:
\begin{enumerate}
    \item a transition algorithm \textbf{Tr} that takes two labellings for the vertices of a connected, static graph $C$ with labels from the set $X$, the graph $C$, and the problem instance, runs in time
at most $f(|C|,x)$, and returns \textbf{true} or \textbf{false},
\item a starting algorithm \textbf{St}, a validity algorithm \textbf{Val}, and a finishing algorithm \textbf{Fin} that all
take a $(k,X)$-state of a connected static graph $C$ and the problem instance, run in time
at most $f(|C|,x)$, and return \textbf{true} or \textbf{false},
\item a vector $\textbf{v}_{upper}$ of $k$ integers, such that $x = (\mathcal{G},\beta)$ is a yes instance of $P$ if and only if there exists a sequence $s_0,...,s_{\Lambda}$ of
$(k,X)$-component states of each snapshot of $\mathcal{G}$ where
\begin{enumerate} 
    \item for each connected component $C_1$ of $G_1$, $\textbf{St}(s_0|_{C_1}, C_1,x) = \textbf{true}$;
\item for each connected component $C_\Lambda$ of $G_\Lambda$, $\textbf{Fin}(s_\Lambda |_{C_\Lambda},C_\Lambda,x) = \textbf{true}$;
\item $\textbf{Tr}(l_{t-1}|_{C_t},l_t|_{C_t},C_t,x) = \textbf{true}$ where $l_i$ is the labelling of vertices of state $s_i$, for all times $1 \leq t \leq \Lambda$ and connected components $C_t$ of $G_t$;
\item $\textbf{Val}(s_t|_{C_t},C_t,x) = \textbf{true}$ for all times $1 \leq t \leq \Lambda$ and connected components $C_t$ of $G_t$; and
\item the sum of vectors satisfies $\Sigma_{0\leq t\leq \Lambda} \Sigma_{C \in C_t} v_{st}(C) \leq \textbf{v}_{upper}$, where $v_{st}$ is the function $v$ in
the $(k,X)$-component state $s_t$ and $\leq$ denotes element-wise vector inequality.
\end{enumerate}
\end{enumerate}
\end{definition}

\begin{lemma} \label{lem:1Qf state}
   The \textsc{temporal role colouring} problem is a $(1,Q,f(C))$-component-exchangeable temporally uniform problem where $f(C)=\rho|C|^3$.
\end{lemma}

\begin{proof}
We want to show \textsc{temporal role colouring} is a $(1,Q,f(C))$-component-exchangeable temporally uniform problem where $f(C)=\rho|C|^3$. That is for $k=1$, $X=Q$ and $f(|C|,(\mathcal{G},A))=\rho|C|^3$ the conditions given in Definition~\ref{def:kxf problem} hold.
We first show that conditions $1$ and $2$ hold for our problem by construction.
First, for labellings $l_i,l_j: V(C) \rightarrow Q$ let \textbf{Tr}$(l_i,l_j,C,(\mathcal{G},A))=\textbf{true}$ if for all $v \in C(V)$ there exists a transition $\delta(l_i(v),c(N_C(v)))=l_j(v)$ and \textbf{false} otherwise. Checking the transition function takes time $\rho$ and finding neighbourhoods and their colours takes time $|C|^2$. Hence the time to check the transition function is $\rho|C|^3$ and Condition $1$ holds.

Next, we define \textbf{St}$(l, C,(\mathcal{G},A))$, as \textbf{true} if, for all $v \in V(C)$, we have $\delta(q_0,\varepsilon)=l_0(v)$, and \textbf{false} otherwise. This runs in time $\rho|C|$. We do not need an extra validity check so we let \textbf{Val}$(l,C,(\mathcal{G},A)) =\textbf{true}$ for all inputs. Finally, we define \textbf{Fin}$(l,C,(\mathcal{G},A))=\textbf{true}$ if for all $v\in C: l(v)\in F$, else \textbf{false}. This takes time $|C|$ to check. Therefore, as these all can be checked in time $\rho|C|^3$, Condition $2$ holds. 

Now we consider Condition $3$. We let $\textbf{v}_{upper}$ be the one dimensional zero vector. We want to show that $x=(\mathcal{G},A)$ is a yes instance if and only if there exists a sequence $s_0,...,s_\Lambda$ of $(1,X)$-component states of each snapshot of $\mathcal{G}$ with the properties stated.

First, we suppose there exists a sequence $s_0,..,s_\Lambda$ of $(1,X)$-component states of $\mathcal{G}$ such that:
\textbf{St}$(s_0|_{C_0},C_0,x)=$\textbf{true} for each connected component $C_0$ of $G_0$, \textbf{Fin}$(s_\Lambda|_{C_\Lambda},c_\Lambda,x)=\textbf{true}$ for each connected component $C_\Lambda$ of $G_\Lambda$ and, for all $1\leq t \leq \Lambda$, \textbf{Tr}$(l_{t-1}|_{C_t},l_t|_{C_t},C_t,x) = \textbf{true}$ for each connected component $C_t$ of $G_t$ .
Then for each vertex $v \in V$ by considering the sequence of connected components such that $v \in V(C_i)$ we have that $\delta(q_0,\varepsilon)=l_0(v)$, $\delta(s_\Lambda(v),N^\Lambda(v))\in F$ and, for all $1\leq t \leq \Lambda$, $\delta(l_{t-1}(v),c(N^t_{C_t}(v)))=l_t(v)$.
As each $C_t$ is a connected component we know that $N^t_{C_t}(v)=N^t(v)$. This gives us that the word  $c(N^1(v))c(N^2(v))...c(N^\Lambda(v))$ is accepted by the automaton. This holds for all $v \in V$ and is exactly the requirement for $(\mathcal{G},A)$ to be a yes instance.



Now we show the converse. Let $(\mathcal{G},A)$ be a yes instance. As \textbf{Val}  is always \textbf{true}, Condition $(d)$ holds vacuously. Similarly, Condition $(e)$ holds as all our vectors are zero. 
As it is a yes instance we have a colouring $r:V \rightarrow [k]$ such that for each vertex $v \in V$ there is an accepting path in $A$ given by $r(N^1(v))r(N^2(v))...r(N^\lambda(v))$  with states $q_1,...q_{\Lambda+1}$ such that $r(v,t)=c(q_t)$. We define the states $s_0,..s_\Lambda$ by $l_i(v)=q_{i+1}(v)$ and show the conditions hold. 
For Condition $(a)$ we know that the only transitions from $q_0$ are $\varepsilon$-transitions, and $q_0$ is our start state of the automaton, therefore $\delta(q_0,\varepsilon)=q_1(v)=l_0(v)$ must be a valid transition for all $v\in V$. That is for all $v\in V$ and therefore for all $V(C_1) \subseteq V$ we have that $\textbf{St}(s_0|_{C_1}, C_1,x) = \textbf{true}$.
For Condition $(b)$ we know that for a word to be accepted by an automaton it must end in an accepting state. In $A$ the accepting states are given by $F$, so we know $l_{\Lambda}(v)=q_{\Lambda+1}(v) \in F$ for all $v \in V$. That is for all $v \in V$ and therefore all $V(C_\Lambda) \subseteq V$ we have that \textbf{Fin}$(s_\Lambda|_{C_\Lambda},c_\Lambda,x)=\textbf{true}$. Finally for Condition $(c)$, we consider a vertex $v$ in connected components $C_1,..,C_\Lambda$. We know that for all $v\in V$ the path $r(N^1(v))r(N^2(v))...r(N^\Lambda(v))$ has states $q_1,...,q_{\Lambda+1}$, which means that each transition from state $q_t$ to state $q_{t+1}$ is labelled by the colours of the neighbourhood $N^t(v)$, that is, for all $1\leq t \leq \Lambda$, $\delta(q_t,r(N^t(v)))=q_{t+1}$. As $r(v,t)=c(q_t)$ and all neighbours of $v$ at time $t$ are in its connected component $C_t$ we can rewrite $r(N^t(v))$ as $c(N_{C_t}^t(v))$. Using this and the $l_i(v)=q_{i+1}(v)$ identity gives us that $\delta(l_{t-1}(v),c(N^t_{C_t}(v)))=l_t(v)$ for all times $1 \leq t \leq \Lambda$. As we showed this for an arbitrary $v\in V$ we have for all times $1\leq t \leq \Lambda$ and connected components $C_t$ that $\textbf{Tr}(l_{t-1}|_{C_t},l_t|_{C_t},C_t,x) = \textbf{true}$.
Therefore we have shown that \textsc{temporal role colouring} is a $(1,Q,\rho|C|^3)$-component-exchangeable temporally uniform problem.
\qed \end{proof}

\end{document}
